\subsection{The exact dyadic arithmetic presentation and its coupled Frobenius words}

\begin{result}{E.185}{Proved}{The pro-two comparison and the one-prime free product}
Let \(S=\{2,\infty,q_1,\ldots,q_t\}\), where the \(q_i\) are
distinct odd primes. Write \(G_S\) for the full Galois group
unramified outside \(S\), and \(G_S(2)\) for its maximal pro-two
quotient. Inflation and local restriction give isomorphisms
\[
 H^2(G_S(2),\mathbb F_2)
   \xrightarrow{\ \sim\ }H^2(G_S,\mathbb F_2)
   \xrightarrow{\ \sim\ }
 H^2(\mathbb R,\mathbb F_2)
       \oplus\bigoplus_{i=1}^t H^2(\mathbb Q_{q_i},\mathbb F_2).
 \tag{E.185.1}
\]
Each summand on the right has order two. There is a minimal
pro-two presentation with generators
\(c,a,x_1,\ldots,x_t\) and defining relations
\[
 c^2=1,\qquad
       y_i x_i y_i^{-1}=x_i^{q_i}\quad(1\leq i\leq t).
 \tag{E.185.2}
\]
Here \(c\) maps to complex conjugation, \(x_i\) to a chosen
generator of tame inertia at \(q_i\), and \(y_i\) is a word
mapping to an actual arithmetic Frobenius at that same prime.
The generator \(a\) can be chosen dual modulo the Frattini
subgroup to \(\kappa_2(2)\), with its other Kummer coordinates
zero. The words \(y_i\) are not independent generators.

For \(t=0\), the group is \(C_2\amalg_{\mathrm{pro}\text{-}2}\mathbb Z_2\).
For \(t=1\) and \(q_1\equiv3\) or \(5\pmod8\), one may
replace \(a\) by the actual Frobenius at \(q_1\), obtaining
\[
 G_{\{2,q_1,\infty\}}(2)
    \simeq C_2\amalg_{\mathrm{pro}\text{-}2}
                \bigl(\mathbb Z_2\rtimes_{q_1}\mathbb Z_2\bigr)
    = C_2\amalg_{\mathrm{pro}\text{-}2}G_{\mathbb Q_{q_1}}(2).
 \tag{E.185.3}
\]
The two indicated factors are the actual real and chosen local
decomposition groups. This result is an arithmetic presentation,
not a simultaneous local section approximation theorem.
\end{result}

\begin{proof}
We first establish the cohomology comparison, including the
maximal-pro-two step. The coefficient module \(\mathbb F_2\)
is literally \(\mu_2\). Since the \(S\)-class group of
\(\mathbb Q\) is trivial, the Kummer localization sequence
identifies \(H^2(G_S,\mu_2)\) with the two-torsion Brauer
classes unramified outside \(S\). The comparison and purity
maps are those of
\citet[II.1.7(b), pp.~151--152; II.2.9, pp.~170--171]{CohMilneADT},
as used in E.155.5. Brauer reciprocity
\citep[II.2.1--2.2, pp.~163--165]{CohMilneADT}
imposes precisely the sum-zero condition on the local invariants.
Omitting the invariant at two therefore gives the second
isomorphism in E.185.1. In particular its dimension is \(t+1\).

For use in proving the first isomorphism, we give explicitly
the small class-number fact needed over cyclotomic fields.
Put \(K_n=\mathbb Q(\zeta_{2^n})\), \(n\geq2\), and let
\(H_n/K_n\) be its Hilbert two-class field. The extension
\(K_n/\mathbb Q(i)\) is cyclic of two-power degree, totally
ramified at the unique dyadic prime, and unramified elsewhere.
The defining maximal unramified abelian two-extension is
characteristic under field automorphisms, so \(H_n/\mathbb Q(i)\)
is a finite Galois two-extension. Let \(P\) be its group and
\(I\) the inertia group of a chosen dyadic prime. Since
\(H_n/K_n\) is unramified, restriction identifies \(I\) with
\(\operatorname{Gal}(K_n/\mathbb Q(i))\), so \(I\) is cyclic.

Let \(J\) be the closed normal subgroup of \(P\) generated
by \(I\). The quotient extension fixed by \(J\) is unramified
at every finite prime of \(\mathbb Q(i)\). If \(P/J\) were
nontrivial, its finite two-group would have a quotient of order
two, giving an unramified quadratic extension of
\(\mathbb Q(i)\). This contradicts the class number one of
\(\mathbb Z[i]\). Thus \(P\) is normally generated by the
cyclic subgroup \(I\). Its Frattini quotient is consequently
generated by one element, so \(P\) is cyclic. Then \(I\) is
normal, and its normal generation implies \(I=P\). Since
\(I\to\operatorname{Gal}(K_n/\mathbb Q(i))\) is an
isomorphism, \(H_n=K_n\). Hence
\[
                    \operatorname{Cl}(K_n)\{2\}=0.
 \tag{E.185.4}
\]
The same holds for every \(S\)-class group of \(K_n\).

It follows from Kummer localization that
\[
       H^2(G_{K_n,S_{K_n}},\mu_2)
                \longrightarrow \operatorname{Br}(K_n)[2]
       \quad\hbox{is injective}.
 \tag{E.185.5}
\]
Indeed the possible kernel is the \(S_{K_n}\)-ideal-class
group modulo two, which is zero by E.185.4. This explicitly
accounts for the ideal-class term rather than inferring
injection merely from vanishing of a Brauer invariant.

Let \(z\in H^2(G_S,\mathbb F_2)\). For sufficiently large
\(n\), restriction to \(K_n\) kills its Brauer image. At two
the local degree is an even power of two, and the real place
becomes complex. At any fixed odd \(q_i\), the local degree
in \(K_n/\mathbb Q\) is the order of \(q_i\) modulo
\(2^n\). These orders tend to infinity: otherwise a fixed
power of the positive integer \(q_i>1\) would equal one in
\(\mathbb Z_2\), hence would equal one as an integer. Thus
each local degree is eventually even. Restriction multiplies
the corresponding Brauer invariant by this local degree.
Equation E.185.5 now shows that the cohomology class itself
restricts to zero on \(G_{K_n,S_{K_n}}\).

Here is the group-theoretic descent that turns this splitting
into the required pro-two comparison. Represent \(z\) by a
continuous central extension
\[
              1\longrightarrow\mathbb F_2
                \longrightarrow E\longrightarrow G_S
                \longrightarrow1.
\]
Let \(G_K=G_{K_n,S_{K_n}}\); it is a normal open subgroup
of two-power index in \(G_S\). A continuous splitting over
\(G_K\) gives a subgroup \(H\) with
\(E_K=\mathbb F_2\times H\). In particular \(H\) is
normal of index two in \(E_K\), not just an arbitrary
subgroup of two-power index in \(E\). Intersect its finitely
many conjugates under \(E/E_K\), and call the intersection
\(N\). Each conjugate is normal of index two in \(E_K\).
Consequently \(E_K/N\) is an elementary abelian two-group,
and \(E/N\) is a finite two-group. Also
\(N\cap\mathbb F_2=1\), because \(N\subset H\).
Thus the original central extension is pulled back from
the finite two-group extension
\[
       1\longrightarrow\mathbb F_2\longrightarrow E/N
           \longrightarrow G_S/\overline N\longrightarrow1,
\]
where \(\overline N\) is the image of \(N\) in \(G_S\).
This proves surjectivity of inflation in E.185.1.
For injectivity, a central extension of \(G_S(2)\) by
\(\mathbb F_2\) is itself pro-two. Any splitting after
pullback to \(G_S\) factors through \(G_S(2)\), by the
universal property of that quotient. It therefore splits
already over \(G_S(2)\). This proves the first isomorphism.

The Kummer group \(H^1(G_S,\mathbb F_2)\) has basis
\[
             \kappa_2(-1),\ \kappa_2(2),\
                     \kappa_2(q_1),\ldots,\kappa_2(q_t).
 \tag{E.185.6}
\]
Complex conjugation pairs only with the first basis vector.
A tame inertia generator at \(q_i\) pairs only with
\(\kappa_2(q_i)\). Choose \(a\) to complete their images
to the dual basis. Burnside's basis criterion gives a
surjection \(F\twoheadrightarrow G_S(2)\), where \(F\)
is free pro-two on \(c,a,x_1,\ldots,x_t\), and its kernel
\(R\) lies in the Frattini subgroup of \(F\).

The local pro-two group at an odd prime has the tame
presentation
\[
 \langle \tau,\sigma\mid\sigma\tau\sigma^{-1}=\tau^q
             \rangle_{\mathrm{pro}\text{-}2}.
 \tag{E.185.7}
\]
Both generators have infinite two-power order. This local
group is realized inside \(G_S(2)\): the global extensions
\(\mathbb Q(\mu_{2^\infty},q^{1/2^\infty})\), ramified
only at \(2,q\), realize the full unramified pro-two tower
and the full tame two-power ramification tower locally at
\(q\). Thus a chosen decomposition group is the full group
in E.185.7. Choose lifts \(y_i\in F\) of its Frobenius
generators. The words
\[
 r_\infty=c^2,\qquad
             r_i=y_i x_i y_i^{-1}x_i^{-q_i}
 \tag{E.185.8}
\]
belong to \(R\).

The five-term sequence for the free presentation identifies
\[
 H^2(G_S(2),\mathbb F_2)
      =\operatorname{Hom}_{\rm cts}
          \bigl(R/\overline{R^2[F,R]},\mathbb F_2\bigr).
 \tag{E.185.9}
\]
Under this identification, restriction to the real factor
is evaluation on \(r_\infty\), and restriction to the
\(q_i\)-decomposition group is evaluation on \(r_i\).
One may check this directly in a central extension: lift
the named local generators and evaluate its single local
relator in the central kernel. Its value is the local
second-cohomology class. Naturality therefore identifies
E.185.1 with the simultaneous evaluations on E.185.8.
Those relators form a basis of the relation module in
E.185.9.

We include the normal-generation step. Let \(R_0\) be
their closed normal closure in \(F\). If
\(R/R_0\ne1\), some finite two-group quotient \(P\) of
\(F/R_0\) has nontrivial image \(D\) of \(R/R_0\).
The equality of the relation module with the span of the
listed relators would give
\(D=D^2[P,D]\). This is impossible for a nontrivial
normal subgroup of a finite two-group: the conjugation
action on the nonzero elementary abelian Frattini
quotient \(D/\Phi(D)\) is a two-group action, and its
coinvariants are nonzero. For example the latter assertion
follows by a central-series induction on the acting finite
two-group, or by applying the augmentation-ideal nilpotence
of its group algebra over \(\mathbb F_2\).
Thus \(R=R_0\), proving the full presentation E.185.2,
not merely its initial relations in an associated graded
group.

When \(t=0\), only \(c^2\) remains. This recovers the
\(\mathbb Q\) case of
\citet[Theorem~2 and proof, author p.~3]{LitGrasJaulent2009}.
For \(t=1\), the evaluation of \(\kappa_2(2)\) on any
arithmetic Frobenius at \(q_1\) is one precisely when
\((2/q_1)=-1\), that is, when \(q_1\equiv3\) or
\(5\pmod8\). Inertia at \(q_1\) does not change this
evaluation. Therefore \(c,\tau_{q_1},\sigma_{q_1}\)
form a Frattini basis, even if \(\sigma_{q_1}\) has
additional \(\kappa_2(-1)\)- or \(\kappa_2(q_1)\)-values.
Use those three elements as generators from the beginning.
The preceding relation argument gives exactly \(c^2\)
and the one tame relation, on disjoint generator sets.
This proves E.185.3 with the stated actual factors.

For several primes, the presentation retains the actual
Frobenius words. For example choose local Frobenius lifts
fixing the compatible roots of their own uniformizer.
Writing \(\epsilon(u/q)\in\mathbb F_2\) for the quadratic
residue character, their first Frattini coordinates are
\[
 y_i\equiv c^{\epsilon(-1/q_i)}
       a^{\epsilon(2/q_i)}
       \prod_{j\ne i}x_j^{\epsilon(q_j/q_i)}
                         \pmod{\Phi(F)}.
 \tag{E.185.10}
\]
These follow by evaluating E.185.6. If all \(q_i\equiv5\pmod8\)
and their pairwise Legendre symbols are one, then each
\(y_i\equiv a\pmod{\Phi(F)}\). This congruence does not
identify the full words. A free product of \(C_2\) and
\(t\) independent local groups would have \(1+2t\)
Frattini generators, whereas the arithmetic group has
\(t+2\). For \(t>1\) the difference is \(t-1\).
Thus the one-prime decomposition supplies no independent
assignment of all the local Frobenius values for several
support primes. E.159 remains a genuine simultaneous
matching condition.
\end{proof}
